\section{Cross-Modal Alignment via Orthogonal Procrustes: Bridging Signals and Symbols in Knowledge Graphs}

\textbf{Authors}: Bryan Alexander Buchorn $\cdot$ Claude Sonnet 4.6 (Research Collab\-orator)  
\textbf{Affili\-ations}: Independent Researcher $\cdot$ Anthropic  
\textbf{Status}: v1.2.0 (Hardened Enterprise)  
\textbf{Date}: March 2026

---

\subsubsection{Abstract}
Knowledge Graphs (KGs) have histo\-rically been limited to symbolic data, creating a repre\-sentational gap between unstr\-uctured physical signals (e.g., sensor telemetry, waveforms) and conceptual entities. We propose the \textbf{Signal Encoder}, a framework for projecting high-dimensional signal features into a symbolic entity embedding space $\mathcal{E}$. By utilizing \textbf{Orthogonal Procrustes Analysis (OPA)} \cite{schonemann1966procrustes, gower2004procrustes} and Singular Value Decomp\-osition (SVD), we learn an optimal rotation matrix $R$ that maps encoded signals to their symbolic count\-erparts while preserving geometric topology. We define two encoding modalities: \textbf{Statistical Encoding} for low-frequency telemetry and \textbf{Spectral Encoding} (Log-FFT) for high-frequency waveforms. Furthermore, we introduce the \textbf{Canonical Basis Anchor} protocol to prevent geometric drift in multi-hop federated reasoning. The v1.2.0 imple\-mentation utilizes \textbf{Namespace Isolation} to prevent semantic collisions between signal and text entities. Our results demonstrate that this alignment enables "Blind Cross-Modal Reasoning" with sub-millisecond latency, providing a critical repre\-sentational bridge for autonomous industrial and scientific AI.

\subsubsection{1. Introd\-uction}
The integration of physical signals into symbolic reasoning systems is a prere\-quisite for advanced autonomous systems. Current approaches often rely on inter\-mediate text descr\-iptions, which introduce significant latency and semantic loss. We demonstrate that direct latent space alignment via Procrustes rotation provides a more efficient and mathe\-matically stable alternative.

\subsubsection{2. Methodology}

\paragraph{2.1 Feature Extraction}
Signals are first transformed into a candidate feature vector $\vec{x} \in \mathbb{R}^d$.
\begin{itemize}
\item \textbf{Statistical}: 16-dimensional vector of moments (mean, variance) and dynamics (velocity, ZCR).
\item \textbf{Spectral}: Magnitude spectrum obtained via FFT, log-scaled and truncated to the target embedding dimension.

\end{itemize}
\paragraph{2.2 Procrustes Alignment}
We solve for the optimal rotation $R$ that minimizes the Frobenius norm between signal points $X$ and symbolic anchors $Y$:
\begin{equation}
R = \arg\min_{\Omega^T\Omega=I} \| \Omega X - Y \|_F
\end{equation}
The solution is derived via SVD of the covariance matrix $M = Y X^T$:
\begin{equation}
M = U \Sigma V^T \implies R = U V^T
\end{equation}

\subsubsection{3. Stability: The Canonical Anchor}
To ensure consistency across federated hops (SPEC\_005), we enforce a protocol where all Signal Encoders align to a designated \textbf{Root Space} $\mathcal{E}_{root}$. This prevents the accum\-ulation of projection noise inherent in nested SVD trans\-formations.

\subsubsection{4. Implem\-entation (v1.2.0)}
The Signal Encoder is implemented as an extension of the \textbf{THALAMUS} pipeline, utilizing \textbf{Namespace Isolation} (\texttt{signal:}) to prevent entity collisions. The projection is a constant-time matrix-vector multi\-plication, suitable for high-velocity streaming envir\-onments.

\subsubsection{5. Conclusion}
Latent space alignment via Orthogonal Procrustes provides a mathe\-matically robust bridge between physical signals and symbolic knowledge. By treating signals as first-class entities, CEREBRUM enables a new class of multi-modal reasoning appli\-cations.

---
\textbf{References}
\begin{enumerate}
\item Schönemann, P. H. (1966). A generalized solution of the orthogonal Procrustes problem. Psycho\-metrika.
\item Gower, J. C., \& Dijkst\-erhuis, G. B. (2004). Procrustes problems. Oxford University Press.
\item Mikolov, T., et al. (2013). Exploiting simil\-arities among languages for machine translation. arXiv preprint.
\item Smith, S. L., et al. (2017). Offline bilingual word vectors, orthogonal trans\-formations and the inverted softmax. ICLR.
\item Conneau, A., et al. (2017). Word Translation without Parallel Data. ICLR.
\item Artetxe, M., et al. (2018). A robust self-learning method for fully unsup\-ervised cross-lingual mappings of word embeddings. ACL.
\item Buchorn, B. A., \& Sonnet, C. (2026). Cross-Modal Signal Projections in CEREBRUM. SPEC\_008.md.

\end{enumerate}